\subsection{获知的理由}

「请用吧。放心，我不会向你收费的。我还不至于愚蠢到对配合我这边的情况空手而来的人做那种事」\\

周被带到了这家怎么看都觉得很高雅的料亭最深处的包间里。刚一落座，饭菜不知为何就端了上来。

到达时看了一眼口袋里的手机，倒的确已经过了正午，吃午饭也很合适，但完全没有想到会以这种形式和小夜共进午餐。\\

「就当是报答你给慧提供一宿一饭的恩情吧。反正谈话肯定会很长，而且我和你谈完之后还有工作呢」\\

既然有能轻而易举地安排这种地方并且毫不吝啬请客的地位，当然会很忙。特意抽出时间来接慧，单凭这一点，周就觉得她对慧的感情比自己想象的要深得多。\\

「……百忙之中抽出时间，非常感谢。那我就恭敬不如从命了」\\

在这里硬要拒绝反而显得失礼，所以周决定真心实意满怀感激地接受小夜的厚意，动起了摆在面前的饭菜。\\

小夜以优美的举止熟练地用餐，这种时候周会将她的身影和真昼重叠在一起，或许是母女的关系吧。虽然这绝对不能对真昼本人说，但在一些诸如举止动作之类的不经意间，两人确实有相似之处。\\

「你盯着我看什么呢」\\

似乎注意到了周的视线，用怀纸轻轻擦拭嘴角的她抛来了一句话。听起来并不像是在责备，倒不如说疑问的成分更重。\\

「是我失礼了，抱歉」

「反正你也是在想我和那孩子长得很像之类的事情吧？」\\

这个人好可怕。周纯粹地感到了一丝恐惧，但如果把这句话说出口就太没礼貌了，于是他拼命咽了下去，用近乎呻吟的声音，并尽量礼貌且积极地将其转换成「您很了解嘛」作为回应。\\

小夜似乎并没有不高兴，她将视线轻巧地转向周，说出了「因为你看起来满脑子都是那孩子的事嘛」这一周无法完全否认的事实。\\

「肚子差不多也填饱了吧，我们进入正题。无论你问什么我都不会生气，但我不保证一定会回答。如果这样可以，那就接受你的提问吧」\\

在周和小夜差不多吃完饭的时候，小夜抛出了周原本的目的——也就是消除疑问的契机。周重新端正了姿势，正面与她相对。

一开始周还在犹豫要问什么样的问题，但首先，也是昨天事件的起因，更是导致曾经的真昼陷入那种境地的理由，问这个大概是最稳妥的。\\

「……去爱没有血缘关系的孩子，而怠慢亲生的女儿，理由是什么？」\\

为什么要把作为父母的一切都奉献给慧，却不给真昼一丝情分。

这也是慧一直抱有的疑问，那个理由究竟是什么呢？\\

「呵呵」

「有什么好笑的吗」

「不？只是觉得你果然很诚实，是那孩子会喜欢的。因为你问得太直接了」

「如果我在这里拐弯抹角地问，您估计会转移话题吧？」

「我是打算拿出诚意来应对的哦？」

「关键是没说一定会回答吧」

「是啊，看来你有在好好听我说话」\\

小夜微笑着的表情，这一点和真昼并不像。\\

小夜似乎对周的反应很满意，对周某种程度上缺乏恭敬的话语毫不在意，她按下代替呼叫铃的桌上按钮，叫来服务员，指示她们收拾吃完的餐具。\\

在这期间，她向服务员叮嘱要清场，并递上了小费，看来是打算进入正题了。\\

餐具收拾完毕，备好了茶水，再次剩下两人独处时，小夜像是要思考一般，将视线微微上扬，稍作沉吟，然后重新面向周。\\

「说来话长，该从哪里说起呢。姑且，就从回答你的问题开始吧。我先说结论，这是优先级的问题。在我心里有比那孩子更优先考虑慧的理由，就是这么回事」

「……您是说，亲生女儿的优先级很低吗？」

「嗯。在我了解一般常识的前提下，是这么回事」\\

被如此干脆地肯定，周的脑海里瞬间一阵发热，但他还不至于失去理性任由这股热血冲头，只是将烦躁和不快感渐渐沉淀在胸口。\\

从她的态度可以看出来，她并不是没有常识，也不是感情用事，只是因为有这样做的必要，所以才这么做了。这会给真昼带来多大的伤害，作为伤害本人的她即使理解，却依然选择了慧。\\

「我签了协议。和慧的父亲——玲」

「协、议」

「没错。履行作为慧的母亲的职责，这样的协议」

「这是说，慧君的父亲也是在理解了会轻视真昼的前提下，向您提出的吗？」

「我想他是理解最终会导致这种结果的哦。在理解的基础上，那个人向我提出了协议，我也答应了」

「为什么……！」

「因为这能给我带来利益，所以我答应了，很简单吧？」\\

这种仿佛完全不把一个人的存在当回事的轻率口吻，再次让周胸口深处翻滚起沸腾的情绪，但如果在这里大声嚷嚷，这次对谈可能就会就此结束。\\

而且，小夜的话语中并没有那种令人作呕的不快感。那确实是非常干脆的说法，这也让周稍微冷静了一些。\\

「您如此追求的利益到底是什么。至少，那是有着哪怕扰乱一个孩子的人生也要去追求的利益吧」

「是啊。我不知道你听说了多少，但那个人和我是作为父母的棋子而结成的婚姻」

「这我知道」\\

因为小夜和朝阳并不是因为相爱而结婚的，所以自己是不被期望出生的孩子——真昼曾这样描述自己，她长大的环境也让她不得不这么说。\\

「哼，看来那孩子相当信任你呢。那她父母的父母的事情呢？」

「哎？」

「对那孩子来说就是祖父母的事情。听说过吗？」

「……没有」

「想必也是。因为和那孩子，可以说是完全没有牵扯」\\

从真昼的口中，几乎没有听过关于祖父母的事。

那是因为真昼本人似乎也和祖父母几乎没有缘分，甚至说过连见都没见过。\\

「你难道不觉得奇怪吗？如果父母不抚养孩子的话，难道祖父母不该接过去吗？父母一方制造了不和的原因，难道不该承担责任吗？」

「……这个，我确实想过。一般来说，如果父母不照顾孩子，这个角色应该轮到祖父母来承担」\\

通常，如果父母放弃抚养，应该会联系亲戚，但真昼说过，祖父母连影子都没有。一般来说，如果女儿生了孩子，祖父母也会关心孙女的情况吧。

对于周说出口的疑问，小夜静静地吐出了一口气。\\

「是啊。没变成那样的理由很简单。是因为我们没让他们参与进来」

「为什么……」

「因为觉得可怜啊」

「哎？」

「看着无聊的过去再次上演，是很让人火大的呢，令人作呕」\\

不同于刚才，这次的话语中带上了感情，且非常强烈。

去年春天看到她时，她对真昼说着严厉的话，周还有些不礼貌地觉得她是个多少有些歇斯底里的人，但实际见面后才发现，她是一个相当冷静、宽容，甚至游刃有余的人。\\

从这样的人口中，突然蹦出了带着强烈厌恶感的话语，周不由自主地看向了小夜的脸，即使感受到了周的视线，她也毫不掩饰那轻蔑的表情。\\

「我啊，恨透了自己的父母」\\

话题突然转变，周的脑海一瞬间有些混乱，但他明白这是在说小夜的父母，也就是刚才提到的祖父母的事。

小夜对自己的父母没有好感的理由，在周的想象范围内，顶多只能想到是因为被强迫政治联姻之类的，但从小夜的表情来看，这并不是仅凭这一点就能解释的，甚至表面隐隐渗出了憎恶与厌恶感。\\

「理由也很老套。把别人的人生当作消耗品来对待，自以为是世界的中心，思想与能力不符的老不死的，我最讨厌这种人了。一想到他们存在于这个世界上就想吐。被时代抛弃的工业废料，那就是我对父母的看法」\\

接连不断蹦出的话语，无论怎么解释，都沾满了负面的感情。\\

「你，被父母深深爱着吧？看一眼就知道了」

「……我不否认」

「我也不否认这个世界上有正直、美好的人存在，但这个世界上，有着比你想象中多得多的，甚至连伪善者的面具都懒得戴的丑恶的人」

「您是想说您的父母就是那样的人吗」

「没错。那些家伙，除了自己以外……无论是孩子还是孙子，都只认为是自己可以利用的棋子，简直是脑子里长了蛆一样的连腐烂的人类都不如的生物。不过对于被小心翼翼养大的你来说，大概是无法理解这种感觉的吧」\\

周，是被父母珍爱着长大的。

虽然因为害羞偶尔会有一些反抗，但打心底里从未讨厌过、疏远过父母，周也尊敬着父母，带着作为家人的爱意活到了现在。至少从周围听到的情况来看，自己大概是在一个相当正直、理想的普通家庭中长大的。\\

所以，周无法想象小夜成长的环境究竟是怎样的。\\

「那种，不该成为父母的人，在这个世界上还是有不少的。从那种父母那里生下来的孩子很悲惨哦，因为一切都会被当作资源来对待」

「资源是指」

「就是资源。可以被使用、可以被消费的消耗品。我没打算博取同情，但我就是从那种父母那里生下来被使用的人，甚至后悔过自己被生下来。你有被强迫牺牲自己去奉献过吗？没有吧」

「……没有」

「想必也是。这是好事。如果这种事情在这个世界上泛滥，那就太让人头疼了」\\

这并不是在迁怒，也不是在讽刺，小夜以一种纯粹觉得那样比较好的表情看着周，轻轻按住自己大臂附近揉搓着。\\

「我啊，恨自己的血。构成这具身体的细胞，这具身体里流淌的鲜血。一想到这些原本是从那些家伙那里挤出来的东西，我就觉得毛骨悚然」\\

「甚至无数次想过自杀」小夜嫌弃地抛下这句话，表情依然严峻。\\

「当然，我也活了这么多年，对自己本身已经妥协了……但怎么可能去爱流着自己血的孩子呢。因为那个笨蛋的愚蠢行为，在我不情愿的情况下怀上了孩子，我怎么可能拥有母性。连对自己都会产生厌恶感的时候，为什么我能去爱，而且还是那个人的孩子呢」

「所以您才对真昼」

「我并没有爱过她哦。虽然觉得她可怜。生在这样的父母家里真是气数已尽了呢。不过我还是安排好了让她能好好活下去的一切」\\

如果厌恶感太强，搞不好甚至会伪装成事故……周也曾想象过这种事，但对小夜来说，或许还没到需要做到那种地步的程度，又或者是对亲自动手没有兴趣，真昼就这样平安地长大了。\\

坚守住了绝对不能跨越的底线的小夜，仿佛要吐出积压在内心的东西一般，长长地叹了一口气。\\

「我特别讨厌那个异端，也讨厌那个笨蛋」

「那个笨蛋……是指朝阳叔叔没错吧」

「嗯。……那个人，连反抗的气魄都没有，只会老老实实地服从，甚至没有自我，是个懦夫。选择忍耐，想着只要等时间过去就会自然消亡的没骨气的懦夫。也不想想快要自然消亡的到底是谁。真是愚蠢」\\

看来小夜对朝阳怀有一种与对父母的厌恶不同的厌恶，她用大拇指和食指触摸着皱起的眉间像是在放松，低着头的同时一瞬间用手掌遮住了脸。

当她再次抬起头时，已经恢复了最初那游刃有余的表情。\\

「我没有道理陪他一起耗下去。别说那些悠哉的话了，排除掉他还比较快呢」

「排除，难道说」

「请别误会。是合法的哦。我并没有对他的身体造成过什么伤害」\\

「请虽然很想伤害就是了」小夜补充了一句危险的话，同时浮现出即使是旁观者也会感到背脊发凉的美丽微笑。\\

「我排除了对任何人来说都是祸害的那些家伙。嘛，与其说是排除，不如说几乎是他们自作自受。不过，为他们铺好通往地狱的路的人毫无疑问是我呢。……对你来说，这不也是一件好事吗？」

「为什么对我来说是好事？」

「如果我没有采取行动，那孩子就会被那些腐烂的老头老太婆当成可以利用的存在了。你甚至应该感谢我才对」\\

啊。\\

周没能忍住发出声音，但小夜并没有责怪他。\\

刚才的话因为是过去的事，周完全没有想象过，但正如小夜所说，如果不是她的行动，「真昼被利用的另一种情况」或许也是有可能发生的。\\

对于小夜提出的这种过于可怕且近在咫尺的可能性，周感到不寒而栗，小夜则像是在调侃一般投来了一句「你满脑子都是那孩子的事呢，真好懂」。\\

「无论那孩子的价值，还是对我来说的价值，我都觉得被认为毫无意义、毫无价值才是更好的情况。如果被认为有利用价值也会很麻烦的，被当成人质就糟糕透顶了。如果只是放任不管，感觉又会被人随便夺走」

「……我觉得有别的说法吧」

「对你来说的价值，和对我来说的价值并不是相等的。……对你来说，这不是一件可以安心的事吗？你最喜欢的孩子已经可以走上自由的人生了。至少不会走上被谁消费的人生。我，也不会去消费那孩子」

「……不消费」

「那就换个说法，不拿父母的权限去支配、去利用那孩子。从一开始我就没有给她强加过任何东西」\\

看着满脸无奈耸着肩膀的小夜，确实，真昼并没有被任何人要求，而是自己主动想要成为父母心目中理想的孩子。哪怕小夜连看都不看她一眼。\\

想要做一个好孩子，是真昼的意志。\\

「看你好像挺不高兴的，觉得我不知道你女朋友有多努力一样」\\

小夜似乎连周之后在想什么都看穿了，她伴随着低语向周露出了微笑。\\

「真好懂呢。看你这副表情，我想你也明白这是那孩子自己选择的路」\\

仿佛一切都被看穿了，周夹杂着惊讶、愤怒和钦佩而保持沉默，小夜依旧是愉悦地笑着你，安静地喝了一口已经完全冷掉的茶。\\

「有用的东西我会用，但我也是会挑的哦。如果做出那种把别人的人生本身都消耗殆尽的举动，我知道是会遭到报应的。因为那些家伙就被他们亲手培养、万分疼爱的部下们背叛，绝望得失去了活下去的力气」\\

小夜若无其事地谈论着明显非常不得了的事情，从中，周依然能隐约感受到她对他们的憎恶。\\

「毕竟他们结下的怨恨多如牛毛，这也是理所当然的结果呢。我和他们都在磨牙霍霍，他们竟然没注意到，真的是老糊涂了吧，呵呵。老了之后理解能力和记忆力都变差了，真是辛苦呢」

「……您的父母，现在」

「正在品尝因果报应之中吧。他们注定要在收容设施里，一边看着一个接一个人掉进地狱，一边在恐惧中悲惨地腐朽。我既然在特等席上欣赏了他们如此难看地没落的样子，自然非常清楚利用别人的人生有多么危险」\\

恐周大概明白了，无论是过去还是将来，真昼绝对不会和祖父母见面。\\

「所以，我只利用我自己。因为是我的人生，我怎么用是我的自由吧？无论结果发生什么，责任由我自己承担。我就是带着这份觉悟，活在当下的」

「……您确实很了不起，但您对真昼所做的事，不就等同于毁了她的人生吗？」

「是啊。所以如果那孩子来复仇，我也觉得无所谓，那也是因果报应。我不说会乖乖接受，但做什么是那孩子的自由」\\

小夜觉得如果真昼想要报复，她会尊重真昼的选择，但真昼恐怕并不期望复仇。单纯地，不想扯上关系的心情应该会更强，而且周也不认为不想惹是生非的真昼会特意去报复。\\

小夜似乎也认定以真昼的性格不会这么做，她像是开玩笑般调侃道「如果她有甩我一巴掌的气魄，我还能在她身上多看出一点可爱之处呢」，随后将视线滑向了周。\\

「说回在刚才的问题里，你觉得最不可理解、最不愉快的协议吧」

「……是指你们两人之间签订的协议吧。我实在难以理解」

「我也没指望你能理解。内容很简单，我来充当慧的好母亲，并给予他适当的教育，而他则要协助我让那些家伙失势。并且要确保以后无论对我们自己，还是对孩子，他们都无法再出手。因为他也是那些人渣的受害者」

「受害者是指」

「那些家伙认为人与人之间的缘分也是可以随便剪贴的东西。因为那些家伙，不只是我，很多人都咽下了苦水，比如说原本预定好的几桩婚事都毁于一旦。净做些多余的事」

「您也是，被卷入其中的一员吧」

「是啊，我本来根本没打算和那个笨蛋成为夫妻的。全都是那些俗物擅自做主」\\

在这一点上，小夜感到愤怒也无可厚非。如果周因为别人的随便安排被迫和真昼分开，被塞给一个完全不认识的人，他也绝对会发火并且去抗议。如果实在没办法，甚至不惜私奔。\\

虽然不知道小夜在被安排给朝阳之前是否有伴侣，但无论有没有，自己的人生计划从根本上被打乱，她无疑会感到愤怒。\\

「我为了建立打垮那些家伙的利益关系，和玲签订了协议。比起我一个人，有个有权力的人在会更顺利。……我事先声明，朝阳也同意了哦。虽然也有无法违抗大势的原因，但毕竟是智惠留下的纪念，他也无法说不吧」

「智惠……？」

「是慧的亲生母亲。虽然慧大概连她的脸都不记得了」\\

因为慧懂事起身边就只有小夜，担负了他过世母亲的职责，也就是说，在他还是婴儿的时候，慧的生母智惠就已经离开人世了。\\

「虽然这不是该由我来说的话，但我觉得那个笨蛋在作为一个人这方面更加有问题。对你来说，或许他看起来大概比我好一点吧」

「您是说，朝阳先生吗？」

「嗯。那个人在你看起来是怎样的？」

「……虽然他对真昼做的事很差劲，但至少我觉得他比您更关心真昼。我觉得他是个温和的人」\\

因为只交谈过一次，所以这只是在谈论一个对不了解的人有什么印象，但周记得他基本上是个能讲得通理的人，散发着温厚且有些虚幻的氛围。\\

「哦。在你看来是这样的啊」

「反过来，在您看来他是怎样的呢？」

「保守地说，是坨屎吧。说是人渣也可以」\\

面对如此堂堂正正、毫不犹豫、断然下定论的小夜，周不禁哑然地看着她的脸。\\

不知道是什么让她如此认为，但既然断言到这个地步并加以辱骂，那两人肯定是非常合不来的。即使考虑到是被父母决定的不情愿的对象，这也未免太刻薄了。\\

「不仅内心脆弱，而且做什么都半途而废，却又自私自利，总是想把好处全占了。虽然这不是该由我来说的话，但他就是个自我中心的人，无论到哪里都是」

「自我中心……是指哪方面？」

「全部呢。对你来说，你以为他是为了那孩子才来谈话的对吧，但那不过是为了他自己而已。那个人其实并不爱任何人。他只爱他自己。无论走到哪里都是如此」\\

面对毫不留情地痛骂朝阳的小夜，周甚至对朝阳产生了一丝怜悯，但小夜仿佛看穿了周的这种想法，用鼻子哼了一声，深深地叹了口气。\\

「你不觉得奇怪吗？」

「奇怪是指哪里？」

「我接下慧母亲的职责，和那个人不照顾那孩子，这两者并不是一回事吧。如果我没法照顾她，那个人难道不应该照顾她吗？再怎么说，也是他在我不情愿的情况下播下种子的罪魁祸首，也是那孩子的父母啊。好歹是父亲，难道不该照顾她吗？至少，玲可是打算哪怕一个人也要照顾慧的」\\

一阵冲击，就好似当头一棒。\\

至今为止，周的思考焦点一直放在小夜放弃抚养上，注意力全被小夜和慧的父亲之间的协议吸引了，但父母可不止一个人。即使小夜因为要负责当慧的母亲而抽不出时间，从她们的话来看，朝阳也是有空闲的。\\

至少，他应该比小夜有更多物理上的空闲去关心真昼。\\

但是，现实却是有一名女孩，被父母双方各自抛弃。\\

「抛弃了她的，是那个人」\\

像是在说教一般，小夜一字一句地用确凿的声音告诉周。\\

「那个人啊，结果只是在逃避，只看自己想看的东西，活在对自己有利的世界里」\\

「那种地方让人毛骨悚然，所以我讨厌他」小夜像看到脏东西一样撂下这句话，周能清楚地感受到她是真的极其厌恶朝阳，虽然不是在说周，但周的胃部还是隐隐作痛。\\

「您所承受的痛苦和这么做的理由我明白了。……尽管如此，作为我来说，就算有理由，也忍不住会想，为什么要对真昼做出那种事」

「为什么能对亲生女儿做到那个地步，是吗？」

「……我觉得普通人是无法做到那种地步，说断就断的」

「是啊，打个比方，你能对一个因为被不情愿的对方强行按倒而怀孕，又被周围人强迫不能堕胎，只好无奈生下孩子的母亲说，让她爱这个孩子吗？」\\

周瞬间屏住了呼吸。\\

恐怕，这是在说小夜自己的事吧。\\

被强迫进行不情愿的婚姻，被要求敞开身体，被周围人强迫生下孩子，在痛苦中分娩。

周实在无法说出那么残酷的话。哪怕是关于真昼的事，他也绝对说不出口。\\

这关系到小夜的尊严，作为一个外人，而且还是男性的周，绝对不能轻率地谈论她被迫品尝屈辱和痛苦的事。如果他这么做了，恐怕小夜不仅会鄙视他，甚至会对他产生杀意吧。\\

「在肚子里孕育，母性就会觉醒，就会爱孩子……这种对某些人来说很方便的母性神话，就别拿出来神圣化了。让人作呕」\\

且不论作为父母这是否正确，小夜无法接受真昼，像讨厌蛇蝎一样讨厌朝阳，在小夜的视角来看，这都是理所当然的事。\\

「那个人连反抗父母的胆量都没有。没有珍惜重要之人的胆量、气量和度量，只是顺从。……就这样，那个人，因为软弱而逃避了。逃进了一个对自己方便的梦里。逃避的结果，就是有了那个孩子」\\

小夜用一种仿佛觉得那已经是发生过的事、是过去的往事，并不觉得有什么特别痛苦的声音解释了真昼的生世，周什么也说不出来。\\

「你说的话在世俗上也是正确的，但我很难接受。所以，我离开了那孩子。作为不看那孩子的补偿，我给了她不愁吃穿的生活、教育，以及自由。所以那孩子现在也没有表现出为钱发愁的样子吧？」\\

真昼虽然没有得到父母的爱，但相应地不用担心金钱方面的问题，并且得到了小雪这个十分理解她的人，正直地长大了。\\

「与其让我这种人待在身边，不如让正直的人去教育她，她会成长得更正直。别说那种令人作呕的甜蜜话语，说什么父母的爱就是一切了。比起浸泡在毒液里长大，那样要干净漂亮得多」\\

小夜的安排，如果揭开盖子来看，在伦理上可能并不正确，但也处于不能断言这绝对错误的范畴之内。\\

如果强迫她强行抚养真昼，那真昼才真的会面临残酷的对待吧。

过度干涉和互不干涉，究竟哪个更好周不知道，但就真昼和小夜之间而言，互不干涉似乎给双方都带来了更好的结果。\\

「你骂我没资格当父母也没关系，事实也确实如此，但你不觉得比起直接造成伤害，这样要好得多吗？我为她准备了不愁金钱的环境，没有强迫她进行强制性的婚姻，还为她安排了毫不吝啬地给予教育并亲切对待她的人。至少，比起被硬塞进父母理想的模子里，被强迫进行消耗身心的奉献，我觉得这是一个好得多的环境」\\

听了刚才的话，周明白小夜所罗列的那些假设都是她曾经遭遇过的，也正因如此，周再也无法否定小夜的话了。\\

小夜对真昼的态度实在无法让人肯定，但环境和情况实在是糟糕到了致命的地步。作为被放弃抚养的当事人，真昼也许被允许对小夜表达不满和抱怨，但作为一个外人的周，在知道了当时小夜所处的立场后，大概没有资格去说什么。\\

看到周紧闭嘴唇，没有表现出要再拿真昼出来责备的举动，小夜微微一笑「你能理解就好」。\\

「一切都太迟了。那孩子觉得我不需要，而且早就脱离了我的掌控。她大概也不想和我有什么和解，对这种事敬谢不敏吧。我也觉得没必要」\\

小夜事到如今也没有打算作为真昼的母亲去对待她，真昼也不把这样的母亲当成父母，拒绝和她扯上关系。双方都选择了拉开距离到无法挽回的地步，互不干涉。

两人之间，一切都可以说是在双方的共识之下，发展到了现在的关系。\\

周做不到再进一步插嘴了。\\

「这样你想问的事情就结束了吧」

「……既然您那么讨厌您的丈夫，为什么，不离婚呢。那样不是更快吗？」\\

周理解了对小夜来说朝阳是该唾弃的存在，但既然如此，为什么不干脆切断缘分呢。\\

「那些下贱家伙的没落花了一段时间，而且在报复的时候，保持婚姻状态更方便。还有，是啊……大概是怜悯吧」

「怜悯？」

「对那个老实、方便、愚蠢的人的同情是一方面。还有，是有人拜托了我。即使是复仇，恩情也是要报答的」\\

小夜用态度表明不打算对周说更多了，周将复仇和恩情这两个词在嘴里咀嚼着。\\

小夜所说的复仇，是对亲生父母，还是对其他人？拜托又是谁拜托的，这都是疑问。拜托她不要离婚，这种请求一般人是不会听的吧。这意味着对方是一个关系好到足以让她听从请求的人。\\

只是，这不是现在该考虑的事，就算问了小夜大概也不会回答。因为这涉及到小夜私人、敏感的部分，所以周想避免再多嘴。\\

「还有一个问题。您为什么，特意去调查真昼的事？听慧君说，这近一年里至少有过一次调查吧」

「啊，连你也一起进行了背景调查让你不开心了？」

「这也是一方面，但从您的立场来看，我原本以为您不会特意在真昼身上花心思的」\\

因为从这次谈话听来，小夜基本上的态度是倾向于完全没有兴趣、没有关心的，所以她特意去打探真昼周围的事情，再次让人感到不自然。\\

「没错。是啊，要是有什么奇怪的东西靠近，利用那孩子，来妨碍我，那不是很困扰吗？」

「您……！」

「对你和那孩子来说又不是什么坏事，你为什么要生气呢。至少，我没有做过对她不利的事」\\

小夜自己应该也知道这话说得很难听，但周觉得她像是在故意用这种挑衅的说法，更进一步说，甚至能感觉到她似乎在享受周表情的变化。

话虽如此，却感觉不到她在愚弄周，只是一直像在守护孩子一样，甚至让人觉得她看着这一切觉得很欣慰。\\

「那么，你打算怎么办？」

「打算怎么办，是指」

「看起来你很喜欢那孩子，大概是打算背负起那孩子的人生吧？」

「……如果是这样，您打算怎么做？」

「没打算怎么做哦？只是没想到会有这么怪的人。你是为了钱？还是看中长相了？那孩子，长得像我和父亲，倒是很适合观赏呢」

「怎么可能！您，把孩子当成什么了……！」\\

这实在不能当做没听见，周吊起眉角正准备抗议而微微起身，却发现小夜正觉得有趣地笑着，于是又把刚要抬起的腰放了下来。\\

（被耍了）\\

看来小夜确实很享受周的反应，看到周察觉到了她的意图，她似乎也觉得很欣慰地放松了脸颊。\\

「你太嫩了呢。遇到这种事你应该冷静地回避才对。稍微被刺激一下就立刻发火。暴露弱点可没好处，小心别人利用你哦，他们会觉得如果要动你，比起动你本人，向你周围的人下手更好」

「……感谢您的忠告，我铭记在心」

「今后你也要继续磨练自己。如果你想保护那孩子的话」\\

那简直就像是，把真昼托付给周一般的话语。\\

「……您能允许我娶真昼吗？」

「有什么不好的」\\

她答应得太干脆、太平静了，害得周露出了今天最呆傻的表情。\\

虽然对小夜来说是个没有兴趣的孩子，但周本以为她会多少有些抵抗，或者说周不行，或者无法接受之类的。

周由于预想被大幅推翻而僵住，小夜则若无其事地继续说道。\\

「因为是与我无关的孩子，那孩子的人生是她自己的，只要那孩子说好，那就没问题吧」\\

话是这么说没错，但能这么顺利地得到许可完全出乎意料，周为了被反对时准备的内心气势眼看着就萎缩了。\\

「你那是什么表情。是觉得我把那孩子当成自己的一部分，可以随便控制的样子？如果是这样的话，那我可是大大地看错你了，我要下调对你的评价咯」

「我不是那个意思……那个，您真的对她没兴趣吗？不是憎恨，只是没兴趣，这么说没错吧？」

「我确实憎恨那流淌的血液，但我并没有把那当作是我自己的延长线。我不是那种自我和他人界限模糊的人，也不是那种想通过孩子来体验人生的愚蠢又自私的人。我是我，那孩子是那孩子，有很明确的界限」\\

聊到这里，周察觉到小夜对自己和对他人都很严格，同时她也是一个能清楚区分自己和他人界限的人吧。\\

虽然是个难以捉摸的人，作为父母也让人觉得不是人，但她并没有致命地缺乏伦理观，倒不如说正因为具备，所以才能做出冷酷的界限划分，可以干脆地做出切割，把女儿的人生就认为是女儿的人生。\\

这对打算包揽真昼人生的周来说是件好事，但同时一切进行得这么顺利，也让他怀疑是不是有什么内情。\\

仿佛看穿了周的想法一般，小夜尖锐地指出了「我不是说过别表现在脸上吗」。\\

「你啊，如果在这里我回答不行，你大概什么都会做吧？字面意思上的，什么都做。你就是那样的人」\\

周由衷地觉得，她真的是一个对周这个人很有理解的人。\\

最坏的情况下，正如之前考虑过的那样，他打算拿出放弃抚养的证据和证言来谈判，但看小夜的样子，那个心理准备似乎也白费了，就周本人而言，算是松了一口气。\\

「仅仅因为是小孩子就过度轻视也不好。像你这样过分正直的孩子，有时候会做出让人意想不到的事情。我也见识过不少了」

「我就当做是夸奖收下了」

「随你喜欢。既然你说需要，那就请便吧。说到底，等成年的时候，无论我许不许可，你们都会自己决定的吧。根本不需要我的同意什么的。所以，只要当事人双方同意，不就好了吗。请自便」

「……那么，我这边就满怀感激地负责照顾她了。就算您说还给您，我也不会还的哦」

「随你便吧？啊，不过，欣赏你为难的表情也许也挺有趣的呢」

「您……」

「呵呵，开玩笑的。我可没那么闲。没闲到去管小孩子的事」

「说不要轻视小孩子的可是您」

「是啊，所以我才没轻视你们，打算最大限度地尊重你们的。还不够吗？」

「够了」\\

这个人果然性格很不讨人喜欢啊，周一边抱着这种可能会反弹给自己的感想，一边带着一半无奈看着小夜，但小夜依然还是老样子，咕噜咕噜地颤动着喉咙，脸上挂着充满了大人从容的微笑。\\

「真是的，太年轻了啊。人心可是很容易改变的。你虽然轻巧地说要带走她，但没想过变心了之后的事情吗？那孩子，可是相当沉重的。她可不会让你逃走的哦？」\\

这大概是，虽然非常拐弯抹角，但也是出于对周的担忧而提出的忠告吧。\\

背负一个人的人生，这意味着不能轻易地半途而废，说不干就不干。对于这种无法挽回的事，你做好心理准备了吗？她大概是在拐弯抹角地质问吧。\\

对周来说，这个问题事到如今已经不需要再问了。\\

「人心很容易改变，这在现实中是无法否认的」\\

周见过很多人心变化的情况。

好的方向也有，坏的方向也有。\\

感情这种东西，容易变迁，是不确定的东西。\\

但是——\\

「如果说您对父母的憎恨一直持续着，那么我爱真昼的心情一直持续下去，这也是有可能的吧」

「哎呀，说得也是。呵呵，这倒是我的盲点」\\

即使有很多东西会改变，但周相信，也有不会改变的东西，有即使改变也会朝着更好方向发展的感情。他有觉悟用一生去证明这一点。\\

面对周的这句话，也许是没有反驳的材料吧，小夜愉快地扬起嘴角，微微眯起眼睛，依旧是流露出了喜悦的感情。\\

「好吧，我就看看你的这份觉悟能持续多久」

「请自便。无论您期不期待，我都会贯穿至死」

「希望对你和那孩子来说，这不会变成大话」\\

小夜不祝他们幸福，也不诅咒他们不幸，只是保持着守望事情发展姿态。周轻轻叹了口气，然后瞥了小夜一眼。\\

最后，周为了问自己一直在意的事，缓缓开口了。\\

「关于慧君，您打算怎么做呢」\\

慧的事，对周来说无论好坏都只是个补充。正题是真昼的事，慧的情况其实和周无关。\\

只是，因为一起度过了时间，周对慧产生了一些共鸣和同情，对于慧今后能否得到他所追求的信息，得到之后又会受到怎样的对待，周果然还是有些在意。\\

一提到慧的名字，她的脸色就变得难看，对她来说，解释情况也许是个令人头疼的问题。\\

「你随便把别人的麻烦事揽上身也要适可而止哦。你能背负的只有那孩子。别看走眼了」

「……您说得对，但是」

「你不用担心，没什么可担心的。回去之后我会把情况告诉慧的。他也到了差不多能理解的年纪了，而且这也是瞒不住的事情。无论是那孩子的事，还是慧的亲生父母的事」\\

刚才周也听小夜说明了她这边的内情，但慧的亲生父亲玲那边的情况，连周都不知道。包括那部分在内，小夜都会告诉慧的吧。\\

「没必要对你这个外人说更多了吧？至少，我没打算害那孩子，也不打算反对你娶她。我想只要有这两个事实，就能作为对你所担忧之事的回答了吧」

「将来慧君这边有什么对这边采取行动的可能性吗？」

「不能断言没有，但那孩子是个聪明的孩子。他有自己判断，不懂就去问别人的判断力和灵活性，也具备理性和良心。真的是个继承了父母优点的孩子。智惠看到了应该会很高兴吧。在奇怪的地方没有像玲真是太好了，要是像那个人，肯定会更腹黑的」\\

周不知道玲的人品，但从小夜的说法来看，他似乎是个相当难对付的人。\\

「你的不安是，今后和那孩子一起生活，会不会受到过去的伤害，对吧？我不打算进行不必要的接触，但你要问能不能相信我，答案肯定是否定的吧。……也是呢，如果你希望的话，做个公证书什么的都可以，你想做到那一步吗？」

「……不用了，谢谢」

「哦」\\

虽然有些不安，但周无论如何也想不到小夜会改变主意来执着于真昼，而且小夜也有她自己最低限度的底线，那就是密切关注着不让真昼陷入不期望的境地。关于这一点，周坦率地感到感激，也觉得可以相信她。\\

过去本身确实应该警戒，但小夜本人并不需要太多的警戒，这就是周的结论。\\

对于周得出的结论，小夜只是简短地点了点头。\\

「那么，到此结束。我不打算进行不必要的接触。我就在高处静静地看着你的觉悟」

「那得收观赏费了」

「哎呀，该付多少钱好呢？一根无名指那么多？」\\

周被她若无其事地说出的话吓得脸颊差点抽搐，但想起小夜的忠告，他告诉自己要冷静，缓缓摇了摇头。小夜大概也看到了他扭曲的脸，但她只是投来像看到玩具一样的眼神。\\

「那就不必了。这种东西就得自己付钱吧」

「我知道」

「……在暗处被人探查真的让人很不开心呢」

「呵呵，是啊。但还是请你勉为其难地接受吧」

「即使没有爱情，也还是有兴趣和关心的啊」

「你最好不要认为，人对别人的感情是可以简单分为正和负两类的」

「我铭记在心」\\

小夜虽然不爱真昼，但也并非完全没有感情。那份感情究竟是什么样的，大概只有小夜自己能命名吧。\\

而且，周肯定再也没有机会去了解它了。\\

「对了，我还有件事要说」\\

就在周准备起身离开时，小夜制止了他的动作。\\

「你啊，为了那孩子大概什么都会做吧，但应该做的事还是看清楚比较好。我不认为你作为代理人站出来，对那孩子来说就是正确的答案……一味地逃避，不堂堂正正地来问我。也不向我明确表达自己的意志。这样真的好吗？」\\

说她没能面对的指责或许是正确的，但偏偏是由造成这一切的本人来说，就算是周也觉得火大，他打破了和缓的气氛，毫不客气地向小夜投去了锐利的视线。\\

「您难道没想过，正是您曾经的所作所为，才在她心里留下了那么深的阴影吗？」

「是啊。……那又怎样？她没有反抗而是在逃避，这是不争的事实。她这一点像那个人，我最讨厌了。不过还没逃进梦里，这点倒比那个人强」

「……那个梦是指」

「就是无法跨越必须跨越的东西，无法咽下必须咽下的东西，所导致的逃避现实。那个人，现在肯定还在做梦吧。虽然偶尔也会醒来，但说不定哪天，就会一直睡下去了呢」\\

虽然小夜用委婉的表达方式在咒骂，但这次的小夜看起来与其说是在生气，不如说更像是在对一个无可救药的人表达着强烈的无奈和怜悯。\\

「干脆因为愤怒而醒过来还比较好，可是他没骨气到连这都做不到，所以大概是不可能的吧。只知道叹息，却没有任何想要自己去想办法的力气，是个只懂得爱惜自己、只知道张着嘴等食吃的、永远的雏鸟」\\

周不明白，小夜在说什么。\\

「呵呵，一副听不懂我在说什么的表情」\\

果然又表现在脸上了吧，周一边在心里发誓下次一定要练好扑克脸，一边点了点头，小夜回了句「想必也是」。\\

「我只是死都讨厌那个人而已。绝对不会再原谅他了」\\

那是无比美丽的微笑，配上平淡的声音，确凿地主张着不再接受周的任何询问。周也悄悄离开了她内心柔软的部分，闭上了嘴唇。
